% =============================================================================
% Chapter: Data Augmentation and Training Data Strategies
% for Pixel-Wise IR--RGB Alignment
% =============================================================================

\chapter{Data Augmentation and Training Data Strategies}
\label{ch:data_augmentation}

Learning-based pixel-wise alignment of infrared and visible imagery demands
dense, accurately labelled training data---specifically, pairs of misaligned
images together with the ground-truth displacement field that maps every pixel
in one modality to its corresponding location in the other.  Collecting such
labels by hand is prohibitively expensive: a human annotator would need to
specify a displacement vector at every pixel, which is impractical for images
of even modest resolution.  Fortunately, when a modest corpus of
\emph{already-aligned} IR--RGB frame pairs is available (e.g.\ a few thousand
frames from a co-registered multi-sensor rig), a rich and virtually
unlimited training set can be manufactured through synthetic misalignment
generation and carefully designed augmentation pipelines.

This chapter presents a comprehensive treatment of the augmentation and
data strategy toolbox for this task.  Section~\ref{sec:synthetic_misalign}
describes how to produce ground-truth displacement fields by applying known
geometric warps to one modality of an aligned pair.
Section~\ref{sec:augmentation_pipelines} details photometric and geometric
augmentation pipelines for paired multi-modal data.
Section~\ref{sec:curriculum} introduces curriculum learning schedules that
stabilise and accelerate convergence.
Section~\ref{sec:limited_data} addresses strategies for maximising
performance when labelled data are scarce.  Finally,
Section~\ref{sec:gt_displacement} formalises the mechanics of ground-truth
dense displacement field generation.


% -----------------------------------------------------------------------------
\section{Synthetic Misalignment Generation}
\label{sec:synthetic_misalign}
% -----------------------------------------------------------------------------

The core insight of synthetic misalignment training is simple yet powerful:
beginning from an aligned pair $(I_{\mathrm{IR}}, I_{\mathrm{RGB}})$, one
applies a known geometric deformation $\mathcal{T}$ to exactly one of the two
images---say $I_{\mathrm{RGB}}$---to obtain a warped image
$\tilde{I}_{\mathrm{RGB}} = I_{\mathrm{RGB}} \circ \mathcal{T}$.  Because
$\mathcal{T}$ is known analytically, the ground-truth displacement field
$\mathbf{d}(\mathbf{x}) = \mathcal{T}(\mathbf{x}) - \mathbf{x}$ is available
at every pixel without any manual annotation.  The network is then trained to
predict $\mathcal{T}^{-1}$ (or equivalently $-\mathbf{d}$) so that applying the
predicted warp to $\tilde{I}_{\mathrm{RGB}}$ recovers alignment with
$I_{\mathrm{IR}}$.  This paradigm has been validated extensively in
single-modality settings such as medical image registration
\cite{balakrishnan2019voxelmorph} and is directly transferable to multi-modal
alignment.

\subsection{Global Parametric Transforms}
\label{subsec:global_transforms}

The simplest class of synthetic misalignments comprises parametric
transformations with a small number of degrees of freedom.

\paragraph{Homographies.}
A projective homography $\mathbf{H} \in \mathbb{R}^{3\times 3}$ has eight
degrees of freedom and can model the dominant geometric relationship between
two planar views.  Random homographies are sampled by perturbing the four
corner points of the image by random offsets drawn from
$\mathcal{U}(-\rho, \rho)$ and computing the resulting $\mathbf{H}$ via the
Direct Linear Transform (DLT) \cite{hartley2003multiple}.  This is the
strategy adopted by HomographyNet \cite{detone2016deep} and has since become
standard practice.  Values of $\rho$ between 16 and 64 pixels (for
$640\times480$ imagery) span a range from nearly imperceptible shifts to
challenging perspective warps.

\paragraph{Affine transforms.}
A six-parameter affine model captures translation, rotation, scale, and shear
while preserving parallelism.  Affine perturbations are appropriate when the
scene content is approximately planar or at sufficient distance from the
cameras.  Parameters are sampled independently: translations from
$\mathcal{U}(-t_{\max}, t_{\max})$, rotation angles from
$\mathcal{U}(-\theta_{\max}, \theta_{\max})$, scale factors from
$\mathcal{U}(1-s_{\max}, 1+s_{\max})$, and shear coefficients from
$\mathcal{U}(-\gamma_{\max}, \gamma_{\max})$ \cite{jaderberg2015spatial}.

\paragraph{Similarity transforms.}
When only translation, rotation, and uniform scaling are expected (a common
scenario for bore-sighted IR--RGB camera pairs with small mechanical
vibrations), a four-parameter similarity transform provides a parsimonious
model that avoids introducing unrealistic shear artefacts.

\subsection{Local Non-Rigid Deformations}
\label{subsec:local_deformations}

Real-world IR--RGB misalignment is rarely a pure global transform.  Lens
distortion differences between the two sensors, parallax from non-planar
scenes, and atmospheric refraction in long-range imaging all introduce
spatially varying misalignment.

\paragraph{Thin-plate spline (TPS) warps.}
Thin-plate splines provide a principled framework for interpolating
displacements specified at a sparse set of control points to a dense field
that minimises bending energy \cite{bookstein1989principal}.  A regular or
random grid of $K$ control points is placed over the image, and each point is
displaced by a random vector drawn from $\mathcal{N}(0, \sigma_{\mathrm{TPS}}^{2})$.
The resulting TPS interpolation yields a smooth, invertible warp when
$\sigma_{\mathrm{TPS}}$ is chosen appropriately (typically 5--20 pixels for
$K = 4\times4$ grids).  TPS warps naturally model the type of local
distortions introduced by lens calibration residuals.

\paragraph{Elastic deformations.}
Following the highly successful augmentation strategy introduced for
biomedical segmentation by Ronneberger et al.\ \cite{ronneberger2015unet},
random dense displacement fields are generated by sampling per-pixel
displacements from $\mathcal{U}(-1,1)$ and smoothing with a Gaussian kernel
of standard deviation $\sigma_{\mathrm{elastic}}$.  The smoothed field is then
scaled by an intensity parameter $\alpha$.  The pair
$(\alpha, \sigma_{\mathrm{elastic}})$ trades off displacement magnitude against
spatial frequency: large $\sigma$ yields gentle, global undulations while
small $\sigma$ produces high-frequency local distortions reminiscent of barrel
and pincushion aberrations.

\paragraph{Radial distortion models.}
To specifically simulate inter-sensor lens distortion differences, the
Brown--Conrady model \cite{brown1966decentering} can be used directly.  Random
radial distortion coefficients $(k_1, k_2, k_3)$ are sampled and applied to
one modality, producing the characteristic barrel or pincushion patterns that
arise when the two sensors have different optical assemblies.

\subsection{Combining Global and Local Perturbations}
\label{subsec:combined_perturbations}

In practice, the most effective training distributions compose global and
local perturbations.  A typical sampling procedure is:
\begin{enumerate}
    \item Sample a global affine or homography $\mathcal{T}_{\mathrm{global}}$.
    \item Sample a local deformation field $\mathcal{T}_{\mathrm{local}}$
          (TPS or elastic).
    \item Compose: $\mathcal{T} = \mathcal{T}_{\mathrm{local}} \circ
          \mathcal{T}_{\mathrm{global}}$.
    \item Compute the dense displacement field
          $\mathbf{d}(\mathbf{x}) = \mathcal{T}(\mathbf{x}) - \mathbf{x}$.
\end{enumerate}
This composition ensures the network encounters a realistic mix of global
shifts (from mechanical vibration or thermal expansion of the camera mount)
and local warps (from lens distortion residuals or parallax), preventing
overfitting to any single deformation mode
\cite{detone2016deep, rocco2017convolutional}.


% -----------------------------------------------------------------------------
\section{Augmentation Pipelines for Paired Multi-Modal Data}
\label{sec:augmentation_pipelines}
% -----------------------------------------------------------------------------

Beyond synthetic misalignment, standard data augmentation remains essential for
regularisation and generalisation, but its application to \emph{paired}
multi-modal data requires care to distinguish augmentations that must be
applied identically to both modalities from those that should differ.

\subsection{Photometric Augmentations}
\label{subsec:photometric}

Photometric augmentations are applied \emph{independently} to each modality,
since the appearance statistics of IR and RGB channels are fundamentally
different.

\begin{itemize}
    \item \textbf{Brightness and contrast.}  Random affine intensity
          transformations $I \mapsto \alpha I + \beta$ with independently
          sampled gains and biases for each modality.  IR imagery often
          exhibits narrower dynamic range, so the sampling ranges should be
          calibrated per-modality.
    \item \textbf{Noise injection.}  Sensor noise characteristics differ
          markedly: visible sensors are dominated by shot noise (approximately
          Poisson) and read noise (Gaussian), while uncooled microbolometer IR
          sensors exhibit significant fixed-pattern noise (FPN) and $1/f$
          temporal noise \cite{nugent2013infrared}.  Augmentation should
          inject Gaussian noise at moderate $\sigma$ for RGB and a combination
          of Gaussian plus column-correlated (stripe) noise for IR to
          faithfully model each sensor's artefacts.
    \item \textbf{Blur.}  Random Gaussian blur with kernel sizes drawn from a
          discrete set (e.g.\ $\{3,5,7\}$) is applied independently.
          Defocus and motion blur kernels provide additional realism.
    \item \textbf{Gamma correction and histogram perturbation.}  Random gamma
          curves $I \mapsto I^{\gamma}$ with $\gamma \in [0.7, 1.4]$ model
          the non-linear response functions common to both sensor types.
\end{itemize}

These photometric augmentations force the alignment network to become invariant
to appearance differences between modalities---a critical capability since
IR--RGB correspondences must be established despite dramatic radiometric
discrepancies \cite{aguilera2016learning}.

\subsection{Geometric Augmentations: Alignment-Preserving vs.\ Misalignment-Creating}
\label{subsec:geometric_augmentations}

A crucial distinction exists between two categories of geometric
augmentation:

\paragraph{Alignment-preserving augmentations.}
Random flips, $90^{\circ}$ rotations, crops, and rescalings applied
\emph{identically} to both $I_{\mathrm{IR}}$ and $I_{\mathrm{RGB}}$ (and to
the displacement field, with appropriate coordinate adjustment) serve as
standard regularisers without altering the alignment relationship.  These
augmentations increase effective dataset diversity while preserving the
ground-truth label.

\paragraph{Misalignment-creating augmentations.}
The synthetic warps of Section~\ref{sec:synthetic_misalign} are applied to
\emph{one modality only} and constitute the misalignment-creating category.
It is essential that these two categories are implemented in separate pipeline
stages to prevent accidental corruption of ground-truth labels.

\subsection{Temporal Augmentations from Video}
\label{subsec:temporal}

When the aligned corpus originates from a video feed, the temporal dimension
provides a natural source of additional training signal.  Consecutive frames
$(I^{t}_{\mathrm{IR}}, I^{t}_{\mathrm{RGB}})$ and
$(I^{t+\delta}_{\mathrm{IR}}, I^{t+\delta}_{\mathrm{RGB}})$ with small
temporal offsets $\delta$ exhibit slight natural misalignments due to
independent sensor readout timing, mechanical vibrations, and scene motion.
These near-aligned pairs can serve as \emph{weakly-labelled} training samples
under the assumption that the true displacement is small.
Additionally, cross-temporal pairs
$(I^{t}_{\mathrm{IR}}, I^{t+\delta}_{\mathrm{RGB}})$ introduce realistic
appearance changes (illumination shift, moving objects) that increase
robustness without requiring explicit photometric augmentation
\cite{liu2020learning}.

\subsection{CutOut, MixUp, and Mosaic for Paired Data}
\label{subsec:cutmix}

Recent augmentation techniques from object detection and classification can
be adapted for paired alignment data with appropriate modifications.

\begin{itemize}
    \item \textbf{CutOut} \cite{devries2017cutout}: Random rectangular regions
          are zeroed in both modalities at \emph{the same spatial location},
          forcing the network to exploit global context rather than relying on
          local texture matches.
    \item \textbf{MixUp} \cite{zhang2018mixup}: Two training triplets
          $(I_{\mathrm{IR}}^{a}, \tilde{I}_{\mathrm{RGB}}^{a},
          \mathbf{d}^{a})$ and
          $(I_{\mathrm{IR}}^{b}, \tilde{I}_{\mathrm{RGB}}^{b},
          \mathbf{d}^{b})$ are linearly blended with weight $\lambda \sim
          \mathrm{Beta}(\alpha, \alpha)$.  The displacement labels are blended
          with the same $\lambda$, which is valid because displacement field
          regression is a linear target space.
    \item \textbf{Mosaic augmentation} \cite{bochkovskiy2020yolov4}: Four
          training pairs are tiled into a $2\times2$ mosaic.  Each quadrant
          retains its own displacement field, and the composite field is simply
          the concatenation of the four sub-fields.  This is particularly
          effective for small datasets as it exposes the network to four
          distinct scenes in a single forward pass.
\end{itemize}


% -----------------------------------------------------------------------------
\section{Curriculum Learning for Alignment}
\label{sec:curriculum}
% -----------------------------------------------------------------------------

Training a dense alignment network directly on the full spectrum of
synthetic misalignments---from sub-pixel shifts to large perspective
warps---can lead to slow convergence and unstable gradients, especially when
using photometric loss components.  Curriculum learning
\cite{bengio2009curriculum} addresses this by structuring the training
schedule from easy to hard examples.

\subsection{Progressive Misalignment Magnitude}
\label{subsec:progressive}

The most natural curriculum for alignment is over the magnitude of synthetic
misalignment.  Training is divided into stages:
\begin{enumerate}
    \item \textbf{Stage 1 (warm-up):} Only small perturbations are applied
          ($\rho \leq 8$\,px for homographies, $\sigma_{\mathrm{TPS}} \leq 3$\,px
          for local warps).  The network learns basic gradient-following
          behaviour and develops stable feature representations.
    \item \textbf{Stage 2 (ramp-up):} Perturbation ranges are linearly or
          exponentially increased over a fixed number of epochs, eventually
          reaching the full target range.
    \item \textbf{Stage 3 (full difficulty):} The complete perturbation
          distribution is used for the remainder of training.
\end{enumerate}
This schedule has been shown to improve final accuracy by 10--15\% compared
to fixed-difficulty training in analogous optical flow tasks
\cite{ilg2017flownet2}.

\subsection{Coarse-to-Fine Training Strategies}
\label{subsec:coarse_to_fine}

Orthogonal to difficulty scheduling, a coarse-to-fine spatial resolution
curriculum mirrors the multi-scale architecture common in flow estimation
networks \cite{sun2018pwcnet, teed2020raft}.  In early training epochs, the
loss is computed only at coarse pyramid levels (e.g.\ $1/8$ or $1/16$
resolution), where the effective receptive field covers large displacements.
As training progresses, finer pyramid levels are activated and contribute to
the loss with increasing weight.  This avoids early corruption of fine-level
features by gradients from large, poorly estimated displacements.

The two curricula---magnitude and resolution---can be combined
multiplicatively: Stage~1 uses small misalignments at coarse resolution,
Stage~2 introduces larger misalignments while activating finer resolutions,
and Stage~3 trains on the full distribution at full resolution.


% -----------------------------------------------------------------------------
\section{Handling Limited Data}
\label{sec:limited_data}
% -----------------------------------------------------------------------------

A few thousand aligned frames, while sufficient to seed synthetic augmentation,
may still lead to overfitting on scene content (e.g.\ the network memorises
specific buildings or landscapes rather than learning modality-invariant
alignment cues).  This section presents complementary strategies for
maximising data efficiency.

\subsection{Self-Supervised and Semi-Supervised Approaches}
\label{subsec:self_supervised}

Self-supervised photometric losses---such as the structural similarity index
(SSIM) or census-transform-based losses
\cite{meister2018unflow, jonschkowski2020matters}---allow the network to train
on \emph{unlabelled} image pairs by penalising appearance discrepancies after
warping.  While direct photometric losses are unreliable across modalities
(IR and RGB may have inverted contrast), modality-invariant representations
can bridge this gap.  For example, one may compute mutual information
\cite{viola1997alignment} or a learned cross-modal feature similarity
\cite{czolbe2021semantic} as the self-supervised objective.  This enables the
use of any available unpaired or loosely-aligned IR--RGB data as additional
training signal.

Semi-supervised training interleaves batches of synthetically labelled pairs
(with ground-truth displacement) and unlabelled pairs (with only the
self-supervised loss), combining the precision of supervised regression with
the coverage of self-supervised learning \cite{lai2017semi_flow}.

\subsection{Exploiting Video Temporal Structure}
\label{subsec:temporal_structure}

Consecutive video frames provide a powerful inductive bias.  If the alignment
network predicts displacement fields $\mathbf{d}^{t}$ and
$\mathbf{d}^{t+1}$ for frames $t$ and $t+1$, a temporal consistency loss
\begin{equation}
    \mathcal{L}_{\mathrm{temporal}} =
    \left\| \mathbf{d}^{t+1} - \mathcal{W}(\mathbf{d}^{t}; \mathbf{f}^{t \to t+1}) \right\|_1
    \label{eq:temporal_loss}
\end{equation}
penalises sudden jumps in the predicted field, where
$\mathcal{W}(\cdot; \mathbf{f})$ denotes warping by an inter-frame optical
flow field $\mathbf{f}$.  This regulariser is particularly effective in
surveillance and autonomous driving settings where alignment should vary
smoothly across time \cite{lai2017semi_flow}.

\subsection{Transfer Learning from Optical Flow Models}
\label{subsec:transfer}

Dense alignment is structurally identical to optical flow estimation: both
predict a per-pixel 2D displacement field.  Pre-trained optical flow models
such as RAFT \cite{teed2020raft}, trained on large synthetic datasets
(FlyingChairs \cite{dosovitskiy2015flownet}, FlyingThings3D
\cite{mayer2016large}, Sintel \cite{butler2012sintel}), provide feature
encoders and correlation layers that transfer well to the IR--RGB alignment
task.  The transfer learning protocol proceeds as follows:
\begin{enumerate}
    \item Initialise the alignment network with RAFT weights pre-trained on
          FlyingChairs $\to$ FlyingThings3D.
    \item Replace or adapt the input stem if the number of input channels
          differs (e.g.\ concatenating a single-channel IR image with a
          three-channel RGB image requires a $4 \to C$ convolutional stem).
    \item Fine-tune on the synthetic IR--RGB alignment dataset with a reduced
          learning rate for the pre-trained encoder and a higher rate for the
          new or adapted layers.
\end{enumerate}
Empirically, RAFT-based initialisation reduces the required fine-tuning data
by a factor of 3--5$\times$ compared to training from scratch, and improves
end-point error on held-out test pairs significantly
\cite{teed2020raft, sun2021autoflow}.

\subsection{Consistency Regularisation}
\label{subsec:consistency}

A forward--backward consistency constraint provides a strong self-supervisory
signal.  Given a pair $(I_{\mathrm{IR}}, I_{\mathrm{RGB}})$, the network
predicts the forward field $\mathbf{d}_{\mathrm{IR} \to \mathrm{RGB}}$ and
the backward field $\mathbf{d}_{\mathrm{RGB} \to \mathrm{IR}}$.  These
should satisfy the invertibility condition
\begin{equation}
    \mathbf{d}_{\mathrm{IR} \to \mathrm{RGB}}(\mathbf{x})
    + \mathbf{d}_{\mathrm{RGB} \to \mathrm{IR}}\!\bigl(\mathbf{x} +
      \mathbf{d}_{\mathrm{IR} \to \mathrm{RGB}}(\mathbf{x})\bigr)
    \approx \mathbf{0}
    \label{eq:fb_consistency}
\end{equation}
at every non-occluded pixel.  The forward--backward consistency loss
\begin{equation}
    \mathcal{L}_{\mathrm{fb}} =
    \frac{1}{|\Omega|}\sum_{\mathbf{x} \in \Omega}
    \left\| \mathbf{d}_{\mathrm{IR} \to \mathrm{RGB}}(\mathbf{x})
    + \mathbf{d}_{\mathrm{RGB} \to \mathrm{IR}}\!\bigl(\mathbf{x} +
      \mathbf{d}_{\mathrm{IR} \to \mathrm{RGB}}(\mathbf{x})\bigr)
    \right\|_1
    \label{eq:fb_loss}
\end{equation}
is differentiable (through bilinear grid sampling) and can be applied to both
labelled and unlabelled pairs.  This regulariser discourages degenerate
solutions and enforces the physically motivated constraint that alignment is a
bijective mapping \cite{meister2018unflow, jonschkowski2020matters}.


% -----------------------------------------------------------------------------
\section{Ground-Truth Displacement Field Generation}
\label{sec:gt_displacement}
% -----------------------------------------------------------------------------

This section formalises the procedure for generating dense ground-truth
displacement maps from synthetic warps, ensuring correctness and numerical
stability.

\subsection{Coordinate Conventions}
\label{subsec:coordinates}

Let the image domain be $\Omega = \{0, \ldots, W\!-\!1\} \times
\{0, \ldots, H\!-\!1\}$ with pixel coordinates
$\mathbf{x} = (x, y)^{\top}$.  A warp $\mathcal{T}: \mathbb{R}^2 \to
\mathbb{R}^2$ maps source coordinates to target coordinates.  The
\emph{displacement field} is defined as
\begin{equation}
    \mathbf{d}(\mathbf{x}) = \mathcal{T}(\mathbf{x}) - \mathbf{x}
    = \bigl( d_x(\mathbf{x}),\; d_y(\mathbf{x}) \bigr)^{\top}.
    \label{eq:displacement}
\end{equation}
For parametric transforms (affine, homography), $\mathcal{T}(\mathbf{x})$ is
computed analytically.  For TPS and elastic warps, $\mathcal{T}$ is defined on
a dense grid.

\subsection{Forward vs.\ Inverse Displacement Fields}
\label{subsec:forward_inverse}

A subtle but critical distinction arises between the \emph{forward} and
\emph{inverse} displacement fields.  Given an aligned pair
$(I_{\mathrm{IR}}, I_{\mathrm{RGB}})$, the warped image is produced by the
inverse (backward) mapping:
\begin{equation}
    \tilde{I}_{\mathrm{RGB}}(\mathbf{x})
    = I_{\mathrm{RGB}}\!\bigl(\mathcal{T}^{-1}(\mathbf{x})\bigr).
    \label{eq:backward_warp}
\end{equation}
The network receives the pair $(I_{\mathrm{IR}}, \tilde{I}_{\mathrm{RGB}})$
and must predict the displacement field that \emph{unwarps}
$\tilde{I}_{\mathrm{RGB}}$ back to alignment.  This target field is:
\begin{equation}
    \mathbf{d}_{\mathrm{gt}}(\mathbf{x})
    = \mathcal{T}^{-1}(\mathbf{x}) - \mathbf{x}.
    \label{eq:gt_field}
\end{equation}
For parametric transforms, $\mathcal{T}^{-1}$ is obtained by matrix inversion
(or iterative refinement for homographies near singular configurations).
For non-parametric warps, the inverse field is computed by scattering the
forward field onto the target grid and filling holes via nearest-neighbour or
bilinear interpolation.

\subsection{Implementation Details}
\label{subsec:implementation}

\paragraph{Grid generation.}
A regular coordinate grid $\mathbf{G} \in \mathbb{R}^{H \times W \times 2}$
is constructed where $\mathbf{G}[i,j] = (j, i)$.  The displaced grid is
$\mathbf{G}' = \mathbf{G} + \mathbf{d}_{\mathrm{gt}}$, and the warped image
is obtained via differentiable bilinear sampling
\cite{jaderberg2015spatial}:
\begin{equation}
    \tilde{I}_{\mathrm{RGB}} = \mathrm{grid\_sample}(I_{\mathrm{RGB}},
    \mathbf{G}').
    \label{eq:grid_sample}
\end{equation}
All modern deep learning frameworks (PyTorch \texttt{F.grid\_sample},
TensorFlow \texttt{tfa.image.dense\_image\_warp}) provide GPU-accelerated
implementations.

\paragraph{Normalisation.}
Displacement fields may be stored in pixel units or normalised to $[-1, 1]$
relative to the image dimensions.  Normalised coordinates are preferred for
training as they make the loss magnitude independent of image resolution:
\begin{equation}
    \hat{d}_x = \frac{2\, d_x}{W - 1}, \quad
    \hat{d}_y = \frac{2\, d_y}{H - 1}.
    \label{eq:normalise}
\end{equation}

\paragraph{Smoothness verification.}
After generating a synthetic displacement field, one should verify that the
Jacobian determinant $\det(\mathbf{J}_{\mathcal{T}})$ is positive everywhere,
ensuring the warp is locally invertible (fold-free).  Fields with negative
Jacobian regions create self-intersecting mappings that are physically
implausible and should be rejected or re-sampled.  For elastic deformations,
this is controlled by choosing sufficiently large $\sigma_{\mathrm{elastic}}$
relative to $\alpha$ \cite{dalca2019unsupervised}.

\paragraph{Multi-scale displacement pyramids.}
For architectures that predict displacement at multiple resolutions (e.g.\
coarse-to-fine networks), the ground-truth field must be downsampled
consistently.  Displacement \emph{magnitudes} scale with resolution: a
displacement of $\mathbf{d}$ at full resolution becomes $\mathbf{d}/2$ at
half resolution.  Bilinear downsampling followed by magnitude scaling
produces the correct multi-scale supervision targets:
\begin{equation}
    \mathbf{d}^{(l)} =
    \frac{1}{2^{l}} \cdot \mathrm{downsample}(\mathbf{d}, 2^l),
    \quad l = 0, 1, \ldots, L.
    \label{eq:multiscale}
\end{equation}

\subsection{Storage and I/O Considerations}
\label{subsec:storage}

Dense displacement fields at full resolution ($H \times W \times 2$, stored as
\texttt{float32}) can be large.  For a $640 \times 480$ image, each field
consumes approximately 2.3\,MB.  For datasets of several thousand pairs with
multiple augmentation variants, on-the-fly generation during training is
strongly preferred over pre-computation.  The warp parameters (e.g.\ the
eight homography offsets plus the TPS control-point displacements) are stored
compactly (a few hundred bytes) and the full displacement field and warped
image are synthesised in the data loader on the GPU.  This approach enables
effectively infinite augmentation diversity with negligible storage overhead.


% -----------------------------------------------------------------------------
\section{Summary}
\label{sec:augmentation_summary}
% -----------------------------------------------------------------------------

The data augmentation and training strategy outlined in this chapter
transforms a modestly sized set of aligned IR--RGB frame pairs into a rich,
diverse, and virtually unbounded training corpus.  Synthetic misalignment
generation (Section~\ref{sec:synthetic_misalign}) provides exact ground-truth
displacement fields at zero annotation cost.  Carefully partitioned
augmentation pipelines (Section~\ref{sec:augmentation_pipelines}) increase
appearance diversity while maintaining label integrity.  Curriculum learning
(Section~\ref{sec:curriculum}) stabilises training and improves convergence.
Limited-data strategies---transfer learning, self-supervision, temporal
consistency, and forward--backward regularisation
(Section~\ref{sec:limited_data})---extract maximum learning signal from every
available frame.  Together, these techniques make it practical to train
high-accuracy dense alignment networks even when only a few thousand aligned
pairs are available.
